\section{Related Systems} \label{sec:related}
Trio \cite{dblp:conf/vldb/agrawalbshnsw06} was an early system for
managing uncertain data, focusing on the representation on different
forms of uncertainty, including probabilities. It
does not support computation of arbitrary marginals, which is at the core
of modern probabilistic databases. Though it adopts a mediator approach,
it is tied to specific and obsolete versions of PostgreSQL (8.2 or 8.3).

Obsolescence is a common problem for systems of this era: the provenance
tracking system Perm \cite{DBLP:conf/icde/GlavicA09} and the
probabilistic datbase system Orion \cite{singh2008orion} are respectively
unmaintained forks of PostgreSQL 8.3 and 8.4; the distributed Orchestra
system
\cite{DBLP:conf/sigmod/GreenKTBIT07}, which provides
an early implementation of provenance semirings, cannot be compiled because
some of its dependencies are on servers that have disappeared; MystiQ
\cite{DBLP:conf/sigmod/BoulosDMMRS05}, implementing safe query
plan evaluation, requires Java 5.0, which reached end-of-service in
2009.

MayBMS
\cite{huang2009maybms}\footnote{\url{https://maybms.sourceforge.net/}} is a probabilistic database system
implementing the general and compact model of U-relations
\cite{DBLP:conf/icde/AntovaJKO08}. Its query evaluator, SPROUT
\cite{DBLP:conf/icde/OlteanuHK09} was designed for efficiency and is in
particular able to exploit the structure of \emph{safe}
queries~\cite{dalvi2013dichotomy}. It was developed as a fork of
PostgreSQL 8.3, which is obsolete and hard to deploy. Some
effort has been made in keeping the system compilable, however, though it
requires using a virtual machine running an older operating system. It
does not support provenance semirings. We
use MayBMS in our experiments as a state-of-the-art system for
probabilistic query evaluation.

GProM~\cite{bahareh2018gprom}\footnote{\url{https://github.com/IITDBGroup/gprom}}
is a middleware layer that adds provenance support to various
database backends (in particular, Oracle and PostgreSQL). It translates
declarative queries with provenance requirements into SQL code, which is
subsequently executed by the backend database system. GProM supports the
capture of provenance for SQL queries, as well as some Datalog queries.
Some features not present in ProvSQL are support for
transactions and PROV-JSON serialization~\cite{provjson}. On the other
hand, in contrast with ProvSQL, it does not support probabilistic query
evaluation, arbitrary semiring provenance, or provenance of
non-monotone queries. As it is modern, actively maintained, and
feature-rich, we use it as a comparison point in
Section~\ref{sec:experiments}.

ProbLog
\cite{DBLP:journals/tplp/FierensBRSGTJR15}\footnote{\url{https://dtai.cs.kuleuven.be/problog/}}
is an actively maintained and easy-to-use probabilistic programming
system, inspired from Prolog and Datalog. SQL queries
and probabilistic databases can be encoded into ProbLog in a
straightforward way. Data can also be stored in a
SQLite database and ProbLog uses knowledge compilation to compute
probabilities of program outputs. However, our experiments showed it does not
scale, as even though the data is stored in a
database, it is loaded in main memory when needed by the query evaluator,
and none of the usual query optimization infrastructure is used.
Indeed, on the simplest query of our benchmark (query~1 of~\qcust, see
Section~\ref{sec:experiments}), ProbLog ran out-of-memory on the smallest
database scale factor used in our experiments (1~GB). On a database 10
times smaller, it did not complete after 5 hours.
